

```latex
\documentclass{article}
\usepackage{graphicx}
\usepackage{amsmath}
\usepackage{enumerate}

\begin{document}

\section{Design Space}
\subsection{Conceptual Design Space}
\begin{itemize}
    \item Refers to the full range of possibilities for addressing identified problems
    \item Infinite in scope
    \item Examples: Digital vs. physical, manual vs. automatic, input/output modalities
\end{itemize}

\subsection{Realistic Design Space}
\begin{itemize}
    \item Mapping of one or two dimensions of the design
    \item Organizes and suggests possibilities
    \item Tangible representations of design ideas
\end{itemize}

\section{Making Design Concrete}
\begin{itemize}
    \item Design process involves lots of cognitive resources, information, and unknowns
    \item Use tangible representations: sketches, prototypes, written notes
    \item Facilitates managing complexity
\end{itemize}

\section{What is a Prototype?}
\begin{itemize}
    \item Limited representation of an envisioned design
    \item Allows interaction and exploration of suitability
    \item Examples: Paper prototyping, screen sketches, storyboards, cardboard mock-ups, Powerpoint slides, videos, software prototypes
\end{itemize}

\section{Why Prototype?}
\begin{itemize}
    \item Evaluation and feedback are central
    \item Stakeholders can interact with prototypes
    \item Enhances team communication
    \item Encourages reflection and design choices
    \item Answers questions and supports decision-making
\end{itemize}

\section{Low-fidelity Prototyping}
\begin{itemize}
    \item Uses media unlike the final form (e.g., paper, cardboard)
    \item Quick, cheap, and easily changeable
    \item Example: Card-based prototypes with index cards
\end{itemize}

\section{High-fidelity Prototyping}
\begin{itemize}
    \item Prototype resembles the final system
    \item User-driven and interactive
    \item No back-end functionality
    \item Common environments: HTML, Flash, Powerpoint
    \item Includes mockups resembling the final interface
\end{itemize}

\section{Advantages and Disadvantages}
\begin{itemize}
    \item \textbf{Low-fidelity}: Cheaper, evaluates concepts, communicates, addresses screen layouts, proof-of-concept
    \item \textbf{High-fidelity}: Functional, defines navigation, user-driven, living specification
\end{itemize}

\section{Disadvantages}
\begin{itemize}
    \item \textbf{Low-fidelity}: Limited error checking, poor coding specs, facilitator-driven, limited utility after requirements, navigation issues
    \item \textbf{High-fidelity}: Expensive, time-consuming, inefficient for PoCs, not effective for requirements gathering
\end{itemize}

\end{document}
```